%!TEX root = ../template.tex
%%%%%%%%%%%%%%%%%%%%%%%%%%%%%%%%%%%%%%%%%%%%%%%%%%%%%%%%%%%%%%%%%%%%
%% abstract-en.tex
%% NOVA thesis document file
%%
%% Abstract in English
%%%%%%%%%%%%%%%%%%%%%%%%%%%%%%%%%%%%%%%%%%%%%%%%%%%%%%%%%%%%%%%%%%%%

\typeout{NT FILE abstract-en.tex}%

One of the main threats to anonymity systems such as Tor is the so-called Sybil attack, in which the attacker impersonates multiple identities for a single entity. To address this threat, I propose Dittor, an anonymous credential scheme that is resistant to Sybil attacks. 

Dittor uses credentials approved by external entities to validate the unique identity of a Tor network node operator anonymously. This credential system does not allow more than one identity for a single node operator, thus mitigating possible Sybil attacks. While previous attempts sought to identify Sybil nodes in the Tor network for mitigation, Dittor solves the problem fundamentally by associating anonymous credentials, which validate the subject's identity, with a pseudonym, which is then used as an identifier in the Tor network to group an entity's nodes into a single connection without revealing its identity. To achieve this end, mechanisms such as Sybil-resilient Anonymous Signatures and Zero-Knowledge Proofs are used to build a Self-Sovereign Identity system, culminating in a Sybil-resilient Anonymous Signatures with Stateless Issuing system.

I propose a nine-month working method, involving research and development phases, with a timeline for simplifying the objectives.

% Abstract keywords
\keywords{
  Tor \and
  Sybil \and
  Decentralized Identifiers \and
  Self-Sovereign Identity
}
